%% ============================================================
%% CHAPTER 9 — Conclusions and Future Work  (budget: ~3,500 words)
%% ============================================================
\chapter{Conclusions and Future Work}
\label{ch:conclusion}

\section{Summary of Contributions}
\label{sec:conc-summary}
% Revisit C1–C5, each in one paragraph, now with the measured numbers
% filled in (parity margin, resource gap, compaction savings, coverage
% jump from L0 to L2 skills, call reduction factor L). Explicitly answer
% each RQ with a yes/no/qualified sentence.

\section{The Thesis Claim, Revisited}
\label{sec:conc-claim}
% One page: in what sense the essence of an agent runtime does not depend
% on an OS — and in what sense it still does (the cloud API is itself a
% service dependency; be honest about the boundary of the claim).

\section{Limitations}
\label{sec:conc-limitations}
% Consolidated: single hardware family (ESP32-S3); simulated environments
% for ground truth (bench subset only); model landscape shifts under the
% measurements; memory-file growth still relies on model discipline;
% cross-user memory isolation (the shared-MEMORY.md leakage) identified
% but not solved here.

\section{Future Work}
\label{sec:conc-future}
% Ranked: (1) MCU-AgentBench — grow the Ch4 task suite into a community
% benchmark with a simulation-only reproduction path; (2) state folding
% (TinyReAct L3+): the model maintains an O(1) state block, the natural
% continuation of Ch6; (3) speculation-to-FSM compilation as the limit of
% Ch8; (4) per-chat memory namespaces and a firmware-enforced memory
% compaction mechanism; (5) resource-aware scheduling: a policy choosing
% among Ch6 levels / Ch8 depth by battery and network state.

\section{Closing Remark}
\label{sec:conc-closing}
% Two or three sentences. The device did not get smarter — the
% architecture got smaller. What that suggests about where agents can
% now live.
